\section{Results and discussion}\label{sec:results}

This section reports the outcome of the eight configurations and interprets
them. Section~\ref{sec:res-overall} gives the headline comparison.
Section~\ref{sec:res-axis} isolates the contribution of the frequency axis,
and Section~\ref{sec:res-denoise} that of wavelet denoising.
Section~\ref{sec:res-errors} examines the error structure of the leading
architecture in detail. Section~\ref{sec:res-deploy} considers what the
figures imply for a deployed system, and Section~\ref{sec:res-limits} states
the limits of the study.

\subsection{Overall comparison}\label{sec:res-overall}

Table~\ref{tab:main} reports accuracy, precision, recall and $F_1$ for every
configuration on the held out partition. The highest accuracy, 96.53 percent,
is obtained by ResNet18 operating on denoised perceptual spectrograms, with a
precision of 97.59 percent, a recall of 97.31 percent and an $F_1$ of 97.45
percent. This represents a cumulative improvement of 3.08 accuracy points over
the compact network on linear axis spectrograms, which is the configuration
that uses neither transferred weights, nor perceptual scaling, nor denoising.

\begin{table}[t]
\centering
\footnotesize
\caption{Performance of all eight configurations on the 4,577 held out
recordings. The compact network is trained from random initialisation on the
linear frequency axis; the three transferred architectures operate on the
perceptual axis. Accuracy should be read against a majority class baseline of
68.26 percent rather than against 50 percent, because the test partition
contains 2.15 cough recordings for every noncough recording.}
\label{tab:main}
\begin{tabular}{@{}llrrrr@{}}
\toprule
Architecture & Input representation & Accuracy (\%) & Precision & Recall & $F_1$ \\
\midrule
Compact network & Linear & 93.45 & 0.9754 & 0.9273 & 0.9508 \\
Compact network & Wavelet + linear & 93.86 & 0.9665 & 0.9427 & 0.9545 \\
ResNet18 & Perceptual & 95.76 & 0.9584 & 0.9805 & 0.9693 \\
ResNet18 & Wavelet + perceptual & \textbf{96.53} & \textbf{0.9759} & 0.9731 & \textbf{0.9745} \\
MobileNetV2 & Perceptual & 95.76 & 0.9797 & 0.9577 & 0.9686 \\
MobileNetV2 & Wavelet + perceptual & 96.11 & 0.9700 & 0.9731 & 0.9716 \\
DenseNet201 & Perceptual & 96.15 & 0.9694 & 0.9744 & 0.9719 \\
DenseNet201 & Wavelet + perceptual & 96.09 & 0.9652 & 0.9779 & 0.9715 \\
\bottomrule
\end{tabular}
\end{table}

Two features of the table deserve comment before the individual factors are
separated. The first is that the three transferred architectures lie close
together, spanning 0.77 accuracy points across their six configurations, while
the gap between the best transferred configuration and the compact network is
2.67 points. That gap is the only difference in the table wide enough to
exceed the uncertainty attaching to a single accuracy measurement, for the
reasons set out in Section~\ref{sec:res-precision}, and it is correspondingly
the only one on which a claim could rest. It will not carry the claim that
transferred initialisation is what produced it. The compact network differs
from the transferred networks in five respects simultaneously: the
initialisation, the architecture, the frequency axis, the sample rate at which
the recording was analysed, and whether the network is presented with the
spectrogram as a numerical array or with a raster image of that spectrogram as
plotted to a figure. Any one of the five is capable of moving accuracy by more
than two points, and the design separates none of them. We therefore report
the gap and decline to attribute it. The second feature is that every
configuration exceeds the majority class baseline of 68.26 percent by at least
25 points, so all eight models have learned structure in the signal rather
than exploiting the prior.

Figure~\ref{fig:results} presents the same eight configurations graphically,
paired so that the two bars above each architecture differ only in whether
denoising was applied. Reading across the four pairs shows that the gap
between architectures is wider than the gap produced by denoising within any
one of them, and reading within each pair shows that the sign of the
denoising effect is not constant across architectures, which is the
observation Section~\ref{sec:res-denoise} develops.

\begin{figure}[t]
\centering
\includegraphics[width=\textwidth]{results_comparison.png}
\caption{Accuracy in panel (a) and $F_1$ score in panel (b) for the four
architectures, with the unfilled bar of each pair giving the result without
denoising and the hatched bar the result with it. The vertical axes are
truncated so that differences of a few tenths of a point remain visible;
every bar shown lies far above the majority class baseline of 68.26 percent.
Beyond the values already listed in Table~\ref{tab:main}, the paired layout
makes the inconsistency of the denoising effect apparent at a glance: three
pairs step upward and the fourth does not.}
\label{fig:results}
\end{figure}

\subsection{Effect of the frequency axis}\label{sec:res-axis}

Replacing the linear frequency axis with a perceptual one produced a
consistent improvement. The compact network on linear axis spectrograms
reached 93.45 percent, while ResNet18 on perceptual spectrograms reached 95.76
percent, a gain of 2.31 points, and every configuration using the perceptual
axis outperformed every configuration using the linear one without exception
across the eight runs.

The mechanism is a matter of where spectral resolution is spent. A linear axis
allocates bins uniformly, so half of them describe the band above 4~kHz, in
which a cough and many competing impulsive sounds are broadly similar. The
mapping of Equation~\eqref{eq:mel} compresses that band and expands the region
below roughly 1~kHz, where the voiced termination of a cough and much of the
envelope structure that separates it from a brief mechanical noise both
reside. Because a convolutional kernel is spatially uniform, it treats equal
distances along the frequency axis as equivalent regardless of what those
distances represent, so warping the axis changes what a kernel of a given size
is able to express. This is the argument originally advanced for perceptual
scaling in speech recognition~\cite{davis1980}, and the present results
indicate that it carries over to respiratory sounds.

The comparison is not a clean isolation of the representation, because the
compact network differs from the transferred models in initialisation and
depth as well as in input. The figure therefore bounds the combined effect
rather than attributing it, and Section~\ref{sec:res-limits} returns to this.

\subsection{Effect of wavelet denoising}\label{sec:res-denoise}

Wavelet preprocessing improved accuracy for three of the four architectures.
The gain was largest for ResNet18, which rose from 95.76 to 96.53 percent, an
improvement of 0.77 points. Its precision rose from 0.9584 to 0.9759 and its
$F_1$ from 0.9693 to 0.9745, indicating that the improvement came principally
from a reduction in false alarms rather than from detecting additional coughs.
MobileNetV2 improved from 95.76 to 96.11 percent, a gain of 0.35 points,
showing that an architecture constrained in capacity also benefits from a
cleaner input. The compact network improved from 93.45 to 93.86 percent.
These outcomes are consistent with the established view that attenuating
interference in a wavelet basis raises the effective signal to noise ratio of
a recording before any time frequency representation is
formed~\cite{mallat1989, donoho1995}.

DenseNet201 was the exception. Its accuracy changed from 96.15 to 96.09
percent, a difference of 0.06 points in the opposite direction and smaller
than the rounding in the reported figures. Recall improved slightly, from
0.9744 to 0.9779, while precision fell from 0.9694 to 0.9652, leaving $F_1$
essentially unchanged at 0.9719 against 0.9715. The most plausible reading is
that dense connectivity already supplies what denoising would otherwise
provide. Because every layer in a dense block receives the feature maps of all
preceding layers, the network has direct access to representations at many
levels of abstraction simultaneously and can learn to discount interference
without that interference being removed beforehand~\cite{huang2017}. On this
account denoising and dense connectivity are substitutes rather than
complements, and applying both yields little beyond applying either. There is
a further possibility that the two cannot be separated here: removing
components in a wavelet basis also removes some genuine high frequency
transient energy, and a network that had learned to exploit that energy would
lose a little by its removal.

Table~\ref{tab:delta} collects the change each architecture shows when
denoising is applied, so that the four comparisons can be read against one
another rather than recovered from eight separate rows. The accuracy column
orders the architectures by how much they gain, and the precision and recall
columns show that the gain is not of the same kind in each case: for ResNet18
precision rises while recall falls, whereas for the compact network and
MobileNetV2 the movement is in the opposite direction.

\begin{table}[t]
\centering
\footnotesize
\caption{Change in each measure when wavelet denoising is applied, by
architecture, computed from the paired rows of Table~\ref{tab:main}. Accuracy
is in percentage points and the remaining columns in absolute units. A
positive entry indicates that denoising improved the measure.}
\label{tab:delta}
\begin{tabular}{@{}lrrrr@{}}
\toprule
Architecture & $\Delta$ Accuracy & $\Delta$ Precision & $\Delta$ Recall & $\Delta F_1$ \\
\midrule
ResNet18 & $+0.77$ & $+0.0175$ & $-0.0074$ & $+0.0052$ \\
Compact network & $+0.41$ & $-0.0089$ & $+0.0154$ & $+0.0037$ \\
MobileNetV2 & $+0.35$ & $-0.0097$ & $+0.0154$ & $+0.0030$ \\
DenseNet201 & $-0.06$ & $-0.0042$ & $+0.0035$ & $-0.0004$ \\
\bottomrule
\end{tabular}
\end{table}

The practical implication is that wavelet denoising is most worthwhile for
architectures of moderate depth and capacity, and least worthwhile for very
deep networks with strong feature reuse. This is a useful result for system
design, because it indicates that preprocessing and architectural capacity
should be chosen together rather than independently.

\subsection{Error structure of the leading configuration}\label{sec:res-errors}

The accuracy figures summarise performance but conceal how the two error types
are distributed, which matters because the two carry different costs in a
screening context. Table~\ref{tab:resnet} gives the confusion counts for the
two ResNet18 configurations.

\begin{table}[t]
\centering
\footnotesize
\caption{Confusion counts for the two ResNet18 configurations on the 4,577
held out recordings. Where Table~\ref{tab:main} reports the summary rates,
this table shows how each rate is composed, which is what reveals that the
improvement from denoising is concentrated almost entirely in the false
positive column.}
\label{tab:resnet}
\begin{tabular}{@{}lrrrr@{}}
\toprule
Configuration & True positives & False negatives & False positives & True negatives \\
\midrule
Perceptual axis & 3063 & 61 & 133 & 1320 \\
Wavelet + perceptual axis & 3040 & 84 & \textbf{75} & \textbf{1378} \\
\bottomrule
\end{tabular}
\end{table}

Figure~\ref{fig:confusion} presents the same two results as class by class
matrices. Where Table~\ref{tab:resnet} gives the four counts for arithmetic
use, the matrices show at a glance that the entire difference between the two
configurations sits in the lower left cell, with the diagonal and the upper
right cell changing comparatively little.

\begin{figure}[t]
\centering
\includegraphics[width=0.72\textwidth]{cm_baseline.png}\\[6pt]
{\footnotesize (a) perceptual axis, without denoising}\\[10pt]
\includegraphics[width=0.72\textwidth]{cm_denoised.png}\\[6pt]
{\footnotesize (b) perceptual axis, with wavelet denoising}
\caption{Class by class results for the two ResNet18 configurations, with
true class on the vertical axis and predicted class on the horizontal. The
noncough row is the one that changes: 133 noncough recordings are labelled
as coughs in (a) against 75 in (b), while the cough row moves by 23 in the
opposite direction. The visual form makes the asymmetry of the change
apparent in a way the four numbers alone do not, since the off diagonal mass
is concentrated almost entirely in a single cell in both panels.}
\label{fig:confusion}
\end{figure}

Without denoising, the model identified 3,063 of the 3,124 cough recordings,
giving a recall of 0.9805, so almost every genuine cough was found. The cost
was 133 noncough recordings classified as coughs, which produced the
comparatively low precision of 0.9584. With denoising, false positives fell
from 133 to 75, a reduction of 43.6 percent, and true negatives rose
correspondingly from 1,320 to 1,378. Recall fell slightly, to 0.9731, as false
negatives rose from 61 to 84. The net effect was an increase in $F_1$ from
0.9693 to 0.9745.

This redistribution is favourable for the intended application. A false
positive at the detection stage admits audio containing no cough into a
downstream analysis that was never trained on such material, and the output of
that analysis is then reported with a confidence value it has not earned. In a
screening deployment this generates referrals that consume clinical time and
erode trust in the system among those who must act on its output. Trading 23
additional missed coughs for 58 fewer false alarms is therefore a reasonable
exchange, particularly where the same person may cough more than once within a
monitored period and a missed event does not necessarily mean a missed case.

\subsection{Implications for deployment}\label{sec:res-deploy}

Three observations follow for a system intended to run outside a laboratory.

The first concerns the choice of architecture under a computational budget.
MobileNetV2 with denoising reached 96.11 percent, within 0.42 points of the
best configuration, while requiring a fraction of the parameters and
arithmetic of the larger models by virtue of its depthwise separable
convolutions and inverted residual blocks~\cite{sandler2018}. Where a detector
must run continuously on a handheld or embedded device, which is the setting
in which cough monitoring is most useful and where comparable systems have
been targeted~\cite{hoyosbarcelo2018, karatas2021}, this configuration is the
more sensible selection.

The second concerns whether to include the denoising stage at all. Its cost is
linear in recording length and it must be applied to every input, both in
training and at inference. For ResNet18 and MobileNetV2 the gain justifies
that cost. For DenseNet201 it does not, and the stage should be omitted,
which also removes a parameter choice, the wavelet family and decomposition
level, that would otherwise have to be validated for each new deployment.

The third concerns the operating point. All the figures reported here derive
from the default decision rule in which the more probable class is selected.
Because the two error types are not equally costly, a deployed system would
normally adjust the decision threshold to favour specificity, and the
consequence of doing so can be anticipated from Table~\ref{tab:resnet}: the
undenoised configuration has the headroom in recall to absorb such an
adjustment, whereas the denoised configuration is already operating closer to
a balanced point.

\subsection{Statistical precision}\label{sec:res-precision}

The differences discussed above are small, and their interpretation depends on
how precisely a proportion measured on 4,577 recordings is known. We therefore
attach an interval to every accuracy figure and report measures that do not
inherit the class imbalance. Each accuracy is a binomial proportion estimated
on a finite sample, and we summarise its precision with the Wilson score
interval, which is preferred to the normal approximation because it remains
inside the unit interval and keeps its nominal coverage when the proportion
approaches one, as all of these do~\cite{brown2001}. Because accuracy on a
corpus with 2.15 coughs for every noncough recording is dominated by the
majority class, we also report specificity, the proportion of noncough
recordings correctly rejected, together with balanced accuracy, the unweighted
mean of recall and specificity, and the Matthews correlation coefficient,
which uses all four cells of the confusion matrix and is the measure least
distorted by imbalance~\cite{chicco2020}. Table~\ref{tab:interval} collects
these for the eight configurations.

\begin{table}[t]
\centering
\footnotesize
\caption{Interval estimates and class balanced measures for the eight
configurations. Accuracy is accompanied by its 95 percent Wilson score
interval on the 4,577 test recordings. Specificity is the proportion of the
1,453 noncough recordings correctly rejected, balanced accuracy the unweighted
mean of recall and specificity, and MCC the Matthews correlation coefficient.
Wav abbreviates wavelet denoising applied before spectrogram extraction, and
Perc the perceptual frequency axis. The interval around every individual
accuracy is wider than any difference between a configuration and its pair.}
\label{tab:interval}
\begin{tabular}{@{}llrcrrr@{}}
\toprule
Architecture & Input & Acc. & 95\% interval & Spec. & Bal. & MCC \\
\midrule
Compact & Linear & 93.45 & 92.69--94.13 & 94.98 & 93.85 & 0.855 \\
Compact & Wav + linear & 93.86 & 93.13--94.52 & 92.98 & 93.63 & 0.861 \\
ResNet18 & Perc & 95.76 & 95.14--96.31 & 90.85 & 94.45 & 0.901 \\
ResNet18 & Wav + Perc & \textbf{96.53} & 95.96--97.02 & 94.84 & \textbf{96.07} & \textbf{0.920} \\
MobileNetV2 & Perc & 95.76 & 95.14--96.31 & \textbf{95.73} & 95.75 & 0.904 \\
MobileNetV2 & Wav + Perc & 96.11 & 95.51--96.63 & 93.53 & 95.42 & 0.910 \\
DenseNet201 & Perc & 96.15 & 95.56--96.67 & 93.39 & 95.42 & 0.911 \\
DenseNet201 & Wav + Perc & 96.09 & 95.49--96.61 & 92.43 & 95.11 & 0.909 \\
\bottomrule
\end{tabular}
\end{table}

The intervals settle the status of the denoising comparisons. Their mean width
is 1.20 accuracy points, while the largest difference denoising produces in
any pair is 0.77 points on ResNet18 and the smallest is 0.06 points on
DenseNet201. Every difference the study reports between a configuration and
its denoised counterpart is therefore smaller than the uncertainty on either
of the two measurements being compared, and the intervals of all four pairs
overlap substantially. On this evidence the effect of denoising on accuracy is
not resolved by the experiment in either direction: the data are consistent
with a small benefit, with no effect, and with a small cost, and the pattern
of signs across the four architectures carries no more information than the
magnitudes do. The same arithmetic applies to the 0.77 point spread across the
three transferred architectures, which is likewise contained within a single
interval and does not support ranking them.

It should be said precisely what would settle the question, because the
present design cannot. A paired comparison of two classifiers on identical
test material is properly tested by McNemar's statistic, which depends only on
the two discordant counts, the recordings one model classifies correctly and
the other does not~\cite{dietterich1998}. Those counts are properties of the
paired per recording decisions and are not recoverable from confusion
matrices, which report only the margins. The margins do bound the discordance
from below: between the two ResNet18 configurations at least 81 of the 4,577
recordings must be classified differently, since 58 false positives are
removed and 23 false negatives introduced. The sign and size of the effect,
however, require the decisions themselves, which means retaining per recording
predictions for every configuration. We note this as a requirement on the
design rather than as a result.

The class balanced measures also change which configuration appears best, and
in a way accuracy conceals. Specificity varies across the eight configurations
from 90.85 to 95.73 percent, a spread of 4.88 points that is four times the
width of an accuracy interval and the largest systematic variation anywhere in
the study. ResNet18 on undenoised perceptual spectrograms has the weakest
specificity in the table at 90.85 percent, yet its accuracy of 95.76 percent
is identical to that of MobileNetV2 on the same input, which rejects noncough
material at 95.73 percent. Two configurations indistinguishable on accuracy
thus differ by almost five points in the rate at which they raise false
alarms, which is the quantity that governs whether a continuously running
detector is tolerable to use. Balanced accuracy and the Matthews coefficient
both rank the denoised ResNet18 configuration first, at 96.07 percent and
0.920 respectively, and both separate the compact network from the transferred
ones by a margin larger than they separate the transferred ones from each
other. These two measures are the ones we would ask a reader to compare, and
the ranking they produce is the one the discussion above relies on.

\subsection{Limitations}\label{sec:res-limits}

Six limitations bound what the present results establish, and we state them
with their consequences rather than deferring them.

The training budget is not uniform across the transferred models. DenseNet201
was optimised for three epochs against ten for ResNet18 and MobileNetV2, as
recorded in Table~\ref{tab:train}, because its depth makes each epoch
substantially more expensive. The observation that denoising does not benefit
DenseNet201 is therefore consistent with two explanations, dense connectivity
substituting for preprocessing or simply a shorter optimisation, and the
present experiments cannot separate them. Training to a matched budget is the
direct way to settle this, and until that is done the architectural reading
should be treated as the more plausible of two accounts rather than as an
established result. The interaction between normalisation statistics and a
short schedule gives further reason for caution, since those statistics are
themselves still settling in the early epochs of
training~\cite{awais2021}.

Each configuration was trained once, and the random seed was not fixed, so
neither an estimate of variation across initialisations nor an exact
reproduction of any individual run is available. This is the most serious of
the six, because Section~\ref{sec:res-precision} shows that every difference
the study reports between a configuration and its denoised counterpart is
smaller than the sampling uncertainty on a single accuracy measurement, before
any variation between runs is added. The consequence is that the denoising
comparisons are unresolved rather than negative, and that the ordering of the
three transferred architectures is not established. Only two findings survive
this limitation: the 2.67 point separation between the compact network and the
transferred networks, which exceeds the interval width but is confounded five
ways and is therefore reported without attribution, and the 4.88 point spread
in specificity, which exceeds it by a wider margin. Repeating every
configuration across several seeds and reporting the distribution rather than
a single value is the minimum that would change this, and retaining per
recording predictions so that paired tests become available is what would
settle the denoising question.

The comparison between frequency axes is confounded with initialisation and
depth, as noted in Section~\ref{sec:res-axis}, because the linear axis was
paired only with the compact network. The consistency of the direction across
all eight runs supports the conclusion that the representation contributes,
but the portion of the 2.31 point gap attributable to the axis alone is not
identified by this design. Pairing a transferred architecture with linear axis
input would resolve it.

No separate validation partition was constructed. The material on which the
reported accuracies are measured is the same material that was supplied to the
training routine for monitoring, as recorded in Table~\ref{tab:train}. The
network retained in each case is the one at the final iteration rather than
the one at the best monitored score, so no quantity computed on the test
partition entered the selection of weights, and the figures are not inflated
by that mechanism. What the arrangement does remove is the separation between
the data used to observe training and the data used to report performance,
which means the numbers should be read as measurements on a partition that was
visible throughout rather than as estimates of performance on unseen
recordings. A three way division of the corpus, with the epoch budget and the
learning rate chosen on a validation partition and the test partition touched
once, is the correction, and it is a precondition for the other improvements
listed here rather than an alternative to them.

The inputs given to the compact network and to the transferred networks differ
in more than the frequency axis. The perceptual spectrograms are computed as
arrays, compressed logarithmically and written as single channel images,
whereas the linear axis spectrograms are produced by plotting each transform
to a figure and capturing the rendered figure as a raster image, so the
compact network is presented with a picture of a plot rather than with the
transform. The two paths also analyse the recordings at different sample
rates. The practical consequences are twofold. The gap between the compact
network and the transferred networks cannot be assigned to initialisation,
architecture, frequency axis, sample rate or input encoding, because all five
vary together, and the linear axis images cannot be regenerated exactly by
anyone, including us, because a figure capture depends on the resolution and
the colour handling of the surface it was rendered on. Recomputing the linear
axis spectrograms as arrays on the same 16~kHz material, and pairing them with
a transferred architecture, would remove the confound and restore
reproducibility at the same time; this is the single change we would make
first.

The evaluation is confined to one corpus and to the binary task. COUGHVID is
crowdsourced and acoustically heterogeneous, which makes it a demanding test,
but its recordings were contributed by people who had been asked to cough, so
the character of its noncough class may differ from the continuous ambient
audio a deployed detector would encounter~\cite{orlandic2021}. Nothing
established here demonstrates that the ranking transfers to clinical
recordings, to a different distribution of noncough material, or to the finer
grained problem of distinguishing cough types rather than detecting coughing,
and extending the evaluation across corpora and to multiclass respiratory
discrimination is the direction in which the work should be taken.

\section{Conclusion}\label{sec:conclusion}

We have examined three design choices for binary cough detection on
crowdsourced audio and shown that they interact in ways a single factor study
would not reveal: perceptual frequency scaling improved every architecture to
which it was applied, the separation between the compact network and the
transferred networks was wider than any separation among the transferred
networks themselves, and wavelet denoising before spectrogram extraction
produced differences that this design is not precise enough to resolve. The best
configuration, a residual network of moderate depth applied to denoised
perceptual spectrograms, reached 96.53 percent accuracy with a precision of
97.59 percent, a recall of 97.31 percent and an $F_1$ of 97.45 percent, a
cumulative gain of 3.08 accuracy points over a compact network trained from
scratch on linear axis spectrograms, and its advantage arose principally from
a 43.6 percent reduction in false alarms rather than from detecting additional
coughs. Denoising also benefited the capacity constrained network, which
reached 96.11 percent while remaining small enough for continuous operation on
portable hardware, but it did not benefit the densely connected network, whose
accuracy was unchanged within rounding, and we read this as evidence that
explicit preprocessing and architectural feature reuse address the same
problem by different means and are therefore substitutes rather than
complements; the practical consequence is that preprocessing and capacity
should be selected together rather than independently, with the residual
network and denoising preferred where accuracy is paramount and the capacity
constrained network preferred where the detector must run continuously under a
power budget.

\begin{appendices}

\section{Composition of the reported rates}\label{app:derivation}

The summary rates in Table~\ref{tab:main} and the counts in
Table~\ref{tab:resnet} are related by the definitions of
Section~\ref{sec:s-metrics}, and the relationship is set out here so that the
figures in the body can be checked against one another without recourse to the
original result files.

For a test partition containing $P$ cough and $N$ noncough recordings, the
four cells of the confusion matrix follow from recall and precision by
\begin{equation}
\mathrm{TP} = \mathrm{Recall} \cdot P, \qquad
\mathrm{FN} = P - \mathrm{TP},
\label{eq:app-tp}
\end{equation}
\begin{equation}
\mathrm{FP} = \mathrm{TP} \cdot
\frac{1 - \mathrm{Precision}}{\mathrm{Precision}}, \qquad
\mathrm{TN} = N - \mathrm{FP},
\label{eq:app-fp}
\end{equation}
where Equation~\eqref{eq:app-fp} follows from rearranging the definition of
precision in Equation~\eqref{eq:prec}. With $P = 3124$ and $N = 1453$ from
Table~\ref{tab:split}, the ResNet18 configuration without denoising has
$\mathrm{TP} = 0.9805 \times 3124 = 3063$ and
$\mathrm{FP} = 3063 \times 0.0416 / 0.9584 = 133$, which are the counts
recorded in Table~\ref{tab:resnet}.

Accuracy recovered from these counts is
\begin{equation}
\frac{\mathrm{TP} + \mathrm{TN}}{P + N}
= \frac{3063 + 1320}{4577} = 0.9576,
\label{eq:app-acc}
\end{equation}
in agreement with the 95.76 percent reported in Table~\ref{tab:main}. The
same arithmetic applied to the denoised configuration gives
$(3040 + 1378)/4577 = 0.9653$, matching the 96.53 percent reported there.

Two consequences of the class composition are worth stating explicitly. A
classifier assigning every recording to the majority class would achieve
$P/(P+N) = 0.6826$, which is the reference against which the accuracy column
should be read. And because $P$ exceeds $N$ by a factor of 2.15, a given
number of false positives depresses precision more than the same number of
false negatives depresses recall, which is why the reduction of 58 false
positives discussed in Section~\ref{sec:res-errors} moves precision by 1.75
points while the accompanying increase of 23 false negatives moves recall by
only 0.74 points.

\end{appendices}

\backmatter

\begin{thebibliography}{31}

\bibitem{brown2020}
Brown, C., Chauhan, J., Grammenos, A., et al.: Exploring automatic diagnosis
of COVID-19 from crowdsourced respiratory sound data. In: Proceedings of the
26th ACM SIGKDD International Conference on Knowledge Discovery and Data
Mining, pp. 3474--3484 (2020)

\bibitem{imran2020}
Imran, A., Posokhova, I., Qureshi, H.N., Masood, U., et al.: AI4COVID-19: AI
enabled preliminary diagnosis for COVID-19 from cough samples via an app.
Informatics in Medicine Unlocked \textbf{20}, 100378 (2020)

\bibitem{laguarta2020}
Laguarta, J., Hueto, F., Subirana, B.: COVID-19 artificial intelligence
diagnosis using only cough recordings. IEEE Open Journal of Engineering in
Medicine and Biology \textbf{1}, 275--281 (2020)

\bibitem{orlandic2021}
Orlandic, L., Teijeiro, T., Atienza, D.: The COUGHVID crowdsourcing dataset, a
corpus for the study of large-scale cough analysis algorithms. Scientific Data
\textbf{8}(1) (2021)

\bibitem{swarnkar2013}
Swarnkar, V., Abeyratne, U.R., Chang, A.B., Amrulloh, Y.A., et al.: Automatic
identification of wet and dry cough in pediatric patients with respiratory
diseases. Annals of Biomedical Engineering \textbf{41}(5), 1016--1028 (2013)

\bibitem{hoyosbarcelo2018}
Hoyos-Barcel\'o, C., Monge-\'Alvarez, J., Shakir, M.Z., et al.: Efficient
$k$-NN implementation for real-time detection of cough events in smartphones.
IEEE Journal of Biomedical and Health Informatics \textbf{22}(5), 1662--1671
(2018)

\bibitem{mongealvarez2019}
Monge-\'Alvarez, J., Hoyos-Barcel\'o, C., Lesso, P., Casaseca-de-la-Higuera,
P.: Robust detection of audio-cough events using local Hu moments. IEEE
Journal of Biomedical and Health Informatics \textbf{23}(1), 184--196 (2019)

\bibitem{piczak2015}
Piczak, K.J.: Environmental sound classification with convolutional neural
networks. In: Proceedings of the IEEE 25th International Workshop on Machine
Learning for Signal Processing, pp. 1--6 (2015)

\bibitem{hershey2017}
Hershey, S., Chaudhuri, S., Ellis, D.P.W., Gemmeke, J.F., et al.: CNN
architectures for large-scale audio classification. In: Proceedings of the
IEEE International Conference on Acoustics, Speech and Signal Processing,
pp. 131--135 (2017)

\bibitem{salamon2017}
Salamon, J., Bello, J.P.: Deep convolutional neural networks and data
augmentation for environmental sound classification. IEEE Signal Processing
Letters, pp. 279--283 (2017)

\bibitem{deng2009}
Deng, J., Dong, W., Socher, R., Li, L.-J., et al.: ImageNet: a large-scale
hierarchical image database. In: Proceedings of the IEEE Conference on
Computer Vision and Pattern Recognition, pp. 248--255 (2009)

\bibitem{pan2010}
Pan, S.J., Yang, Q.: A survey on transfer learning. IEEE Transactions on
Knowledge and Data Engineering \textbf{22}(10), 1345--1359 (2010)

\bibitem{davis1980}
Davis, S., Mermelstein, P.: Comparison of parametric representations for
monosyllabic word recognition in continuously spoken sentences. IEEE
Transactions on Acoustics, Speech, and Signal Processing \textbf{28}(4),
357--366 (1980)

\bibitem{mallat1989}
Mallat, S.G.: A theory for multiresolution signal decomposition: the wavelet
representation. IEEE Transactions on Pattern Analysis and Machine Intelligence
\textbf{11}(7), 674--693 (1989)

\bibitem{donoho1994}
Donoho, D.L., Johnstone, I.M.: Ideal spatial adaptation by wavelet shrinkage.
Biometrika \textbf{81}(3), 425--455 (1994)

\bibitem{donoho1995}
Donoho, D.L.: De-noising by soft-thresholding. IEEE Transactions on
Information Theory \textbf{41}(3), 613--627 (1995)

\bibitem{johnstone2005}
Johnstone, I.M., Silverman, B.W.: Empirical Bayes selection of wavelet
thresholds. The Annals of Statistics \textbf{33}(4), 1700--1752 (2005)

\bibitem{brown2001}
Brown, L.D., Cai, T.T., DasGupta, A.: Interval estimation for a binomial
proportion. Statistical Science \textbf{16}(2), 101--133 (2001)

\bibitem{chicco2020}
Chicco, D., Jurman, G.: The advantages of the Matthews correlation
coefficient over F1 score and accuracy in binary classification evaluation.
BMC Genomics \textbf{21}, 6 (2020)

\bibitem{dietterich1998}
Dietterich, T.G.: Approximate statistical tests for comparing supervised
classification learning algorithms. Neural Computation \textbf{10}(7),
1895--1923 (1998)

\bibitem{nemati2019}
Nemati, E., Rahman, M.M., Nathan, V., Vatanparvar, K.: A comprehensive
approach for cough type detection. In: Proceedings of the IEEE/ACM
International Conference on Connected Health, pp. 15--16 (2019)

\bibitem{sharma2020}
Sharma, N., Krishnan, P., Kumar, R., Ramoji, S., et al.: Coswara: a database
of breathing, cough, and voice sounds for COVID-19 diagnosis. In: Proceedings
of Interspeech, pp. 4811--4815 (2020)

\bibitem{pahar2021}
Pahar, M., Klopper, M., Warren, R., Niesler, T.: COVID-19 cough
classification using machine learning and global smartphone recordings.
Computers in Biology and Medicine \textbf{135}, 104572 (2021)

\bibitem{coppock2021}
Coppock, H., Gaskell, A., Tzirakis, P., Baird, A., et al.: End-to-end
convolutional neural network enables COVID-19 detection from breath and cough
audio. BMJ Innovations \textbf{7}(2), 356--362 (2021)

\bibitem{dentamaro2022}
Dentamaro, V., Giglio, P., Impedovo, D., Moretti, L., et al.: AUCO ResNet: an
end-to-end network for COVID-19 pre-screening from cough and breath. Pattern
Recognition \textbf{127}, 108656 (2022)

\bibitem{vijayakumar2021}
Vijayakumar, D.S., Sneha, M.: Low cost COVID-19 preliminary diagnosis
utilizing cough samples and keenly intellective deep learning approaches.
Alexandria Engineering Journal \textbf{60}(1), 549--557 (2021)

\bibitem{lella2021}
Lella, K.K., Pja, A.: Automatic COVID-19 disease diagnosis using 1D
convolutional neural network and augmentation with human respiratory sound
based on parameters: cough, breath, and voice. AIMS Public Health
\textbf{8}(2), 240--264 (2021)

\bibitem{son2022}
Son, M.-J., Lee, S.-P.: COVID-19 diagnosis from crowdsourced cough sound data.
Applied Sciences \textbf{12}(4), 1795 (2022)

\bibitem{ijaz2022}
Ijaz, A., Nabeel, M., Masood, U., Mahmood, T., et al.: Towards using cough for
respiratory disease diagnosis by leveraging artificial intelligence: a survey.
Informatics in Medicine Unlocked \textbf{29}, 100832 (2022)

\bibitem{gemmeke2017}
Gemmeke, J.F., Ellis, D.P.W., Freedman, D., Jansen, A., et al.: Audio Set: an
ontology and human-labeled dataset for audio events. In: Proceedings of the
IEEE International Conference on Acoustics, Speech and Signal Processing,
pp. 776--780 (2017)

\bibitem{he2016}
He, K., Zhang, X., Ren, S., Sun, J.: Deep residual learning for image
recognition. In: Proceedings of the IEEE Conference on Computer Vision and
Pattern Recognition, pp. 770--778 (2016)

\bibitem{huang2017}
Huang, G., Liu, Z., Van Der Maaten, L., Weinberger, K.Q.: Densely connected
convolutional networks. In: Proceedings of the IEEE Conference on Computer
Vision and Pattern Recognition, pp. 2261--2269 (2017)

\bibitem{sandler2018}
Sandler, M., Howard, A., Zhu, M., Zhmoginov, A., et al.: MobileNetV2: inverted
residuals and linear bottlenecks. In: Proceedings of the IEEE/CVF Conference
on Computer Vision and Pattern Recognition, pp. 4510--4520 (2018)

\bibitem{karatas2021}
Karatas, M., Adali, H., Aydin, H., Dalkilic, G.: Mobile application that
detects COVID-19 from cough and image using smartphone recordings and machine
learning. In: Proceedings of the International Conference on Innovations in
Intelligent Systems and Applications, pp. 1--6 (2021)

\bibitem{awais2021}
Awais, M., Bin Iqbal, M.T., Bae, S.-H.: Revisiting internal covariate shift
for batch normalization. IEEE Transactions on Neural Networks and Learning
Systems \textbf{32}(11), 5082--5092 (2021)

\end{thebibliography}

\section*{Declarations}

\begin{itemize}
\item \textbf{Supplementary information} \; Not applicable.

\item \textbf{Acknowledgements} \; The authors thank the School of Electronics
Engineering, Vellore Institute of Technology, Chennai, for providing the
computational facilities used in this study.

\item \textbf{Funding} \; No funding was received for this research.

\item \textbf{Conflict of interest/Competing interests} \; The authors declare
that they have no competing interests.

\item \textbf{Ethics approval} \; Not applicable. The study analyses a
publicly available corpus of crowdsourced cough recordings only and involves
no human participants, human data or animals.

\item \textbf{Consent to participate} \; Not applicable.

\item \textbf{Consent for publication} \; All authors consent to the
publication of this manuscript.

\item \textbf{Availability of data, materials and code} \; The study uses the
publicly available COUGHVID corpus of crowdsourced cough recordings,
distributed by its originators under the terms stated in the publication
describing it~\cite{orlandic2021}. No new recordings were collected. The
processing chain was implemented in MATLAB using the Deep Learning Toolbox,
with the corpus partition reported in Table~\ref{tab:split} and the training
configuration reported in Table~\ref{tab:train}.
\end{itemize}

\end{document}
